%% kapa: exam=1 round=32 year=2021 period=1 form=A subject=civil_law subject_ko=민법 seq=7 number=7
%% source: Q-Net 감정평가사 기출문제 (cst00309) artlSeq=5208397 file=extracted/1008/5208397/A형/2021년도 제32회 감정평가사 1차 1교시 A형.pdf
%% conversion: kordoc 4.24.51 -> tools/gen1.py (auto); review=proofread-sonnet; flags=none
\begin{question}{7}
비법인사단에 관한 설명으로 옳은 것은? (다툼이 있으면 판례에 따름)
\begin{choices}
  \item 비법인사단의 대표자는 자신의 업무를 타인에게 포괄적으로 위임할 수 있다.
  \item 여성은 종중구성원이 되지만, 종중총회의 소집권을 가지는 연고항존자가 될 수는 없다.
  \item 이사의 선임에 관한 민법 제63조는 비법인사단에 유추적용될 수 없다.
  \item 교회는 비법인사단이므로 그 합병과 분열이 인정된다.
  \item 비법인사단의 대표자가 총회의 결의를 거치지 않고 총유물을 권한 없이 처분한 경우에는 권한을 넘은 표현대리에 관한 민법 제126조가 준용되지 않는다.
\end{choices}
\end{question}
\answer{5}
